\begin{annexitem}{annex:p19-nl-ratio-identification}{thm:p19-nl-ratio-identification}{Comparison equivalence with different mean codes}
\begin{proof}
The random variable \(e^{\varepsilon/\kappa}\) is positive.
Cauchy--Schwarz and the assumed finite exponential second moment
give
\(\mathbb E(e^{\varepsilon/\kappa})
 \leq\sqrt{\mathbb E(e^{2\varepsilon/\kappa})}<\infty\).
Its expectation is strictly positive: the positive events
\(\{e^{\varepsilon/\kappa}\geq1/j\}\), for positive integers
\(j\), cover the sample space, so at least one has positive
probability and supplies a positive expectation lower bound.
Thus division by \(m\) is well defined.
The log code's mean is \(\kappa\ln x\); adding a constant does
not change variance, so its variance is
\(\operatorname{Var}(\varepsilon)\), independently of \(x\).
By definition \(\mathbb E(Z_x)=1\), and
\(\mathbb E(Z_x^2)=\mathbb E(e^{2\varepsilon/\kappa})/m^2\)
is finite. Multiplication by \(x>0\) therefore gives
\(\mathbb E(S_x)=x\),
\(\operatorname{Var}(S_x)=x^2\operatorname{Var}(Z_x)\)
and the asserted standard deviation, including the zero-variance case.

For each draw, positivity and strict monotonicity of the exponential
give the equivalences
\[
 \begin{split}
 L_x>L_y
 &\ \Longleftrightarrow\
   \ln x+\varepsilon_x/\kappa>
     \ln y+\varepsilon_y/\kappa\\
 &\ \Longleftrightarrow\
   xe^{\varepsilon_x/\kappa}>ye^{\varepsilon_y/\kappa}\\
 &\ \Longleftrightarrow\ S_x>S_y.
 \end{split}
\]
Equality is also preserved, so no unmentioned assumption about
ties is needed. Independently identically distributed noise makes
the distribution of \(\varepsilon_y-\varepsilon_x\) the same
for every pair. The first comparison event is precisely
\(\{\varepsilon_y-\varepsilon_x<\kappa\ln(x/y)\}\);
its probability therefore depends only on the positive ratio.
Coupling the two constructions with the same noise yields identical
comparison outcomes for every pair while their stipulated means
have different forms. Observing those comparisons alone cannot
identify which mean/noise decomposition generated them.
No claim of equivalence for recall, estimation locations, response
times or neural recordings follows from this calculation.
\end{proof}
\end{annexitem}

\begin{annexitem}{annex:p19-nl-proportion-bias}{thm:p19-nl-proportion-bias}{Bias patterns and coincident response mappings}
\begin{proof}
For \(\beta>0\), the functions \(t^\beta\) and
\((1-t)^\beta\) are continuous on \([0,1]\), interpreting
\(0^\beta=0\). They cannot both vanish at the same point.
Their positive sum therefore makes the quotient continuous, including
at both endpoints. Direct substitution gives
\(P_\beta(0)=0\), \(P_\beta(1)=1\) and
\(P_\beta(1/2)=1/2\). Interchanging the two summands proves
\(P_\beta(1-t)=1-P_\beta(t)\).

For \(0<t<1\), put \(r=t/(1-t)>0\). Division of numerator
and denominator by \((1-t)^\beta\) gives
\[
 P_\beta(t)=\frac{r^\beta}{1+r^\beta},\qquad
 \frac{P_\beta(t)}{1-P_\beta(t)}=r^\beta.
\]
As \(t\) increases in the open interval, \(r\) strictly
increases; raising a positive number to a positive power and applying
\(z\mapsto z/(1+z)\) both strictly preserve order. Hence
\(P_\beta\) strictly increases there, and its interior values
lie strictly between its endpoint values.
Since \(t=r/(1+r)\), comparison of \(P_\beta(t)\) with
\(t\) is equivalent to comparison of \(r^\beta\) with \(r\).
If \(0<t<1/2\), then \(0<r<1\) and \(\ln r<0\).
For \(0<\beta<1\), \(\beta\ln r>\ln r\), so
\(r^\beta>r\) and \(P_\beta(t)>t\). For \(\beta>1\)
the inequality reverses. If \(1/2<t<1\), then \(r>1\)
and \(\ln r>0\); the inequalities reverse relative to the
lower half. At \(\beta=1\) the quotient simplifies to \(t\)
throughout, including the endpoints. This proves all one-cycle cases.

For the two-cycle formula, both definitions give \(1/2\) at
\(t=1/2\), using \(P_\beta(1)=1\) and
\(P_\beta(0)=0\). Each piece is continuous and strictly increasing,
so the matching values join them to a continuous strictly increasing
mapping. Substitution gives fixed values at
\(0,1/4,1/2,3/4,1\).
On the lower half,
\[
 Q_\beta(t)-t=\tfrac12\bigl(P_\beta(2t)-2t\bigr).
\]
On the upper half, with \(s=2t-1\),
\[
 Q_\beta(t)-t=\tfrac12\bigl(P_\beta(s)-s\bigr).
\]
The four quarter intervals send \(2t\) or \(s\) to the
successive lower or upper open halves of \([0,1]\). Applying the
one-cycle cases proves over, under, over, under for
\(0<\beta<1\), with the reversed cases when \(\beta>1\).
For \(\beta=1\), both expressions vanish, giving
\(Q_1(t)=t=P_1(t)\).

Finally, multiplication by \(U>0\) and substitution
\(t=x/U\) gives \(\widehat y=x\) from either mapping when
\(\beta=1\). Every accurately placed target in the stipulated
range therefore yields the same observable prediction from the
two constructions. Such positions alone cannot decide which reference
construction produced them. This is an identification limitation of
these stated mathematical response models; it neither proves that a
participant used either one nor asserts equivalence on other tasks.
\end{proof}
\end{annexitem}
